% Figure: the water analogy, honestly framed (ch01)
\begin{tikzpicture}[scale=0.88,font=\small]
  % ---------- left: water ----------
  \begin{scope}
    % pump
    \draw[thick] (0,0) circle (0.5);
    \node at (0,0) {\scriptsize pump};
    % pipe loop
    \draw[very thick,accent] (0,0.5) -- (0,1.3) -- (4.2,1.3) -- (4.2,0.2);
    \draw[very thick,accent] (0,-0.5) -- (0,-1.3) -- (4.2,-1.3) -- (4.2,-0.2);
    % constriction
    \draw[very thick,accent] (4.2,0.2) -- (4.2,-0.2);
    \draw[thick,fill=graybg] (4.2,0.2) -- (4.55,0.08) -- (4.55,-0.08) -- (4.2,-0.2) -- cycle;
    % flow arrows
    \draw[->,thick] (1.4,1.05) -- (2.2,1.05);
    \draw[->,thick] (2.8,-1.05) -- (2.0,-1.05);
    \node[align=left,anchor=west] at (5.0,0.55) {\scriptsize pump $\approx$ source};
    \node[align=left,anchor=west] at (5.0,0.1) {\scriptsize pressure $\approx$ voltage};
    \node[align=left,anchor=west] at (5.0,-0.35) {\scriptsize constriction $\approx$ resistor};
    \node[align=left,anchor=west] at (5.0,-0.8) {\scriptsize flow $\approx$ current};
    \node at (2.3,-2.0) {\scriptsize\itshape the analogy};
  \end{scope}
  % ---------- right: where it breaks ----------
  \begin{scope}[yshift=-3.4cm]
    \node[align=left,anchor=north west,text width=13.2cm] at (-0.4,0)
      {\scriptsize\textbf{Where it breaks:} cut a pipe and the water pours out; cut a wire and
       charge simply stops. There is no ``leak'' of charge, nothing pushes sideways, and the
       pressure exists \emph{between} two points, not inside one point. Keep the picture for
       intuition in this chapter; drop it the moment it argues back.};
  \end{scope}
\end{tikzpicture}
